\section{Purpose and scope}
\label{sec:introduction}

The intended package supports ordinary Lean proofs about languages with
binders. A researcher should be able to specify syntax, define operations, and
reason about judgments using certified renaming and freshness principles.
The mathematical setting is classical Lean, with one atom sort, actual
quotients, and eventual least supports. Single binders and nested binders form
the initial binding scope. Ordinary propositions remain the language of
theorems; the package is not intended to require a separate object language
for logic.

Before syntax generation or binder reasoning can be certified, the meaning of
renaming must be fixed. A carrier alone does not determine an action: the same
type can carry a trivial action, a canonical atom action, or, for suitable
carriers, a pointwise or conjugation action. These choices affect every later
support statement. The foundation described here therefore exposes a small
permutation and action interface for ordinary Lean proofs.
The underlying mathematics is standard group-action and nominal-set theory;
Pitts~\cite[Chapters 1--2]{pitts2013} provides the background.

For each definition or structure that encodes a mathematical concept, we
state that concept before describing its Lean representation. We identify
the correspondence between the mathematical object, its operations and
laws, and the actual declarations. We also check the pinned Mathlib for
existing constructions before adding package-specific machinery, recording
what is reused and why a different interface is needed when reuse does not
fit. Representation helpers without a separate mathematical counterpart
are described as such, rather than promoted to additional theory.

Declarations use the namespace \texttt{NominalPackage}, under
\path{Package/}, with \path{Package.lean} as
the public import. It uses Mathlib directly and preserves the earlier
\path{Nominal/} and \path{Instances/} development as a reference implementation.
Analogous results in that development are not evidence that a new package
declaration has been proved. The design adapts the subgroup construction for
finite permutations and uses standard \texttt{MulAction} interfaces for the
new boundary.

\paragraph{Implementation status.}
This version describes the implemented foundation increment for roadmap tasks
\texttt{PKG-F02} and the first slice, \texttt{F03a}, of \texttt{PKG-F03}.
The permutation and action interfaces exported by \texttt{Package} have been
checked by Lean, and the source-coverage check and production axiom audit pass.
Section~\ref{sec:evidence} records the verification evidence separately from
the mathematical exposition.

The support-facing slice \texttt{F03b} remains open under \texttt{PKG-F03}.
Supported predicates (\texttt{PKG-01}) are also later work. Persistent usage
examples are reserved for the future case studies.

The present scope includes finite permutations, swaps, actions on atoms,
products and finite atom sets, explicit discrete data, and equivariance of
ordinary functions. It does not establish finite-swap factorization, finite
or least support, freshness, quotient actions, supported predicates, name
abstraction, induction, recursion, or generated syntax. These are later
mathematical layers, with their own proof obligations.
